\documentclass[11pt]{article}

\usepackage{amsmath,amssymb,amsthm}
\usepackage[margin=1.1in]{geometry}
\usepackage{bm}
\usepackage{booktabs}
\usepackage{array}
\usepackage[colorlinks=true,linkcolor=blue,citecolor=blue]{hyperref}

\newtheorem{proposition}{Proposition}
\newtheorem{claim}{Claim}
\newtheorem{remark}{Remark}
\newcommand{\R}{\mathbb{R}}
\newcommand{\eps}{\varepsilon}
\newcommand{\bu}{\mathbf{u}}
\newcommand{\bw}{\mathbf{w}}
\newcommand{\bn}{\mathbf{n}}
\newcommand{\dS}{\,\mathrm{d}S}
\newcommand{\dV}{\,\mathrm{d}V}
\newcommand{\tr}{\operatorname{tr}}
\newcommand{\Div}{\operatorname{div}}
\newcommand{\code}[1]{\texttt{#1}}
\newcommand{\jump}[1]{[\![#1]\!]}

\title{The static response of a body with a stratified fluid core:\\
       well-posedness, the Dahlen closure, and what dissipation\\
       selects}
\author{AdGIA method notes}
\date{7--8 October 2026}

\begin{document}
\maketitle
\sloppy

\noindent
This note collects, in one place and with terms defined, a set of
results and observations about the \emph{static} (time-independent)
linearised response of a self-gravitating body containing an inviscid
fluid region --- a fluid core --- whose stratification is not neutral.
The short version:

\begin{enumerate}
\item If the fluid's squared buoyancy frequency $N^2$ is nonzero, the
  static problem posed with the fluid as an elastic material (the
  natural reading of ``set the time derivatives to zero'') has
  \emph{no solution}: its field equations force the fluid's radial
  displacement to ride the equipotentials, and that assignment is
  generically incompatible with the kinematic matching at the
  core--mantle boundary (\S\ref{sec:slaving}--\ref{sec:nonexist}).
\item This failure is invariant under the fluid-relabelling symmetry:
  no choice of gauge representative repairs it. The static problem is
  well posed on the quotient by the relabelling group \emph{exactly}
  when $N^2 = 0$, where the group is large enough to absorb the
  offending degree of freedom (\S\ref{sec:gauge}).
\item Dahlen's classical treatment of the fluid is immune because it
  eliminates the fluid displacement before the conflict can arise. The
  price, usually left implicit, is a normal-displacement discontinuity
  at the core--mantle boundary, carried in the formalism as a surface
  mass density. Dahlen's static problem is well posed and its answer
  is independent of the fluid's bulk modulus at every degree
  $\ell \ge 1$ (\S\ref{sec:dahlen}).
\item Dissipation decides. Every physical regularisation we can
  identify --- fluid viscosity, diffusion of the stratifying quantity
  (a point made by Wunsch in the 1970s\footnote{Reference to be added;
  the observation that the static response of a stratified fluid is
  selected by dissipative processes, however weak, goes back at least
  to Wunsch's work of that decade.}), or a Maxwell viscoelastic core
  --- yields a well-posed evolution whose $t \to \infty$ limit is the
  Dahlen (secular) answer. Dahlen's closure is therefore not one
  arbitrary resolution among many: it is the distinguished limit of
  every dissipative completion (\S\ref{sec:regular}).
\item For glacial-isostatic-adjustment timescales the body sits
  extraordinarily deep in the secular regime --- even a core with
  $N^2$ six orders of magnitude below plausible values has buoyancy
  periods of years against loading periods of millennia --- so the
  Dahlen closure is the physically correct static limit there, and the
  entire ambiguity quantified below is practically negligible for GIA
  next to, say, three-dimensional viscosity variations. The interest
  of the result is conceptual, and for bodies or observables where
  core stratification itself is the target (\S\ref{sec:persp}).
\item The dissipative selection is not only the theory but now the
  \emph{method}: making the fluid regions an artificial Maxwell
  solid, applying the load as a step, and relaxing until the
  \emph{solid} stops moving computes the secular response with
  nothing but welded elastic solves --- no fluid machinery, no slip
  interface, no kernel handling beyond rigid modes, on any geometry.
  Implemented generically in the problem layer and, after
  optimisation, \emph{cheaper} in three dimensions than the welded
  static solve it replaces (\S\ref{sec:maxwell}).
\item This inverts the ranking of the formulations. Free slip is the
  $\eta \to 0$ asymptotics of a viscous boundary layer --- correct
  about the interface once the layer has collapsed, silent about the
  interior physics that collapses it --- and the measurements show
  the slip classes wandering in the same $N^2$ band as welded: both
  are finite idealisations of the dissipative completion. The
  relaxation is the formulation-faithful computation; Dahlen is its
  converged spherical closure; welded and slip are its
  approximants, retained as instruments (\S\ref{sec:persp}).
\end{enumerate}

\noindent
Numerical evidence for each statement, obtained with this library's
one-dimensional radial reference and its finite-element solvers, is
collected in \S\ref{sec:evidence}. Statements are tagged
\emph{derived} (a calculation reproduced here), \emph{argued} (a
mathematical claim whose proof is sketched but not completed), or
\emph{measured}.

%=====================================================================
\section{Setting and definitions}
\label{sec:setting}

\paragraph{Background state.} A non-rotating body occupies
$B \subset \R^d$ ($d = 2$ or $3$; the two-dimensional ``gravitating
disc'' is the testing geometry of this library and all statements
below hold in both dimensions). It consists of a solid region and a
fluid region separated by an interface $\Sigma$ (for concreteness, a
fluid core of radius $c$ under a solid mantle; the core--mantle
boundary, CMB). The background is a radial hydrostatic equilibrium:
density $\rho(r)$, gravitational potential $\Phi_0$ with
$\nabla^2\Phi_0 = 4\pi G \rho$ and $g(r) = \Phi_0'(r) \ge 0$, and
pressure $p^0$ with $\mathrm{d}p^0/\mathrm{d}r = -\rho g$.

\paragraph{The fluid, and $N^2$.} The fluid is \emph{inviscid} and
\emph{barotropic-elastic}: its strain energy per unit referential
volume is $V(\mathbf{x}, J)$, a function of position and of the
volume ratio $J = \det F$ alone, with bulk modulus
$\kappa = J^2 \partial^2_J V$ evaluated at equilibrium. It supports
no shear stress. The \emph{squared buoyancy frequency} is
\begin{equation}
  N^2 \;=\; -g\left(\frac{1}{\rho}\frac{\mathrm{d}\rho}{\mathrm{d}r}
            + \frac{g\rho}{\kappa}\right),
  \label{eq:N2}
\end{equation}
the frequency (squared) of vertical oscillation of a displaced fluid
parcel. $N^2 > 0$ is stable stratification, $N^2 < 0$ unstable, and
$N^2 = 0$ --- the \emph{Adams--Williamson} or \emph{neutral} state ---
means the density profile is exactly the one produced by adiabatic
self-compression, so a displaced parcel is always neutrally buoyant.
The library's test model \code{fluid\_core} (uniform density,
finite $\kappa$) has $N^2 = -g^2\rho/\kappa < 0$; its constructed
twin \code{aw\_core} has $N^2 = 0$ identically, by choosing the
density profile first and defining
$\kappa = -\rho^2 g / (\mathrm{d}\rho/\mathrm{d}r)$ in closed form.

\paragraph{Perturbations.} Linearise about the background:
$\bu$ is the (referential) displacement field, $\phi$ the
\emph{Eulerian} perturbation of the gravitational potential (the
change at a fixed spatial point), and, in the fluid,
$p^1$ and $\rho^1$ the Eulerian pressure and density perturbations.
The \emph{Lagrangian} pressure perturbation (the change following the
material) is $\Delta p = p^1 + \bu\cdot\nabla p^0$; for the
barotropic-elastic fluid $\Delta p = -\kappa\,\Div\bu$. Mass
conservation gives $\rho^1 = -\Div(\rho\,\bu)$.

\paragraph{The three formulations compared.} All act on the same
background and load; they differ in how the fluid and the interface
are described. (Precise weak forms:
\code{doc/quasi\_static\_models.tex} and
\code{doc/gravitating\_elasticity.md}; here only what matters.)
\begin{description}
\item[welded:] the fluid is kept as the elastic material above, and
  the displacement is a single continuous field across $\Sigma$ ---
  both normal and tangential components match.
\item[slip (slipping interface):] the fluid displacement is an
  independent field; only the \emph{normal} component is matched at
  $\Sigma$ ($\bn\cdot\jump{\bu} = 0$, in the notation of
  \code{doc/slip\_interface.tex}); the tangential jump --- the slip
  --- is free, and the interface carries the corresponding pressure
  and gravity forms.
\item[Dahlen:] the fluid displacement is eliminated analytically
  (\S\ref{sec:dahlen}); the fluid contributes a volume term in $\phi$
  and interface terms at $\Sigma$, and the solid is described in the
  standard (effective-moduli) variables.
\end{description}

\paragraph{The fluid-relabelling (FRL) symmetry.} A
\emph{relabelling} is an infinitesimal rearrangement $\bw$ of the
fluid's material labels that leaves every physical field unchanged:
it must preserve the mass distribution and whatever else the energy
depends on, and must not move the boundary,
$\bw\cdot\bn = 0$ on $\partial(\text{fluid})$. Which $\bw$ qualify
depends on the stratification:
\begin{itemize}
\item \textbf{strict class} ($N^2 \ne 0$): $\Div\bw = 0$ and
  $\bw\cdot\nabla\Phi_0 = 0$ --- volume-preserving and tangent to the
  level surfaces of the background potential. On a spherical
  (circular) interface the induced boundary motions are the
  surface-divergence-free tangential fields; in two dimensions,
  rigid rotations only.
\item \textbf{class K} ($N^2 = 0$): $\Div(\rho\,\bw) = 0$ with
  $\bw\cdot\bn = 0$ --- the ``trivial displacements''. This is a much
  larger group; in particular every tangential boundary field is the
  trace of some member, so every slip is removable
  (\code{doc/planning/open\_issues.md}, the removability item).
\end{itemize}
Solutions of the static equations, when they exist, are determined
only up to the relevant class: the physically meaningful object is
the equivalence class (the \emph{orbit}), and any gauge-invariant
observable --- surface displacement, the potential, tractions --- is
constant on it. The Maitra--Al-Attar programme implemented in this
library enforces normal continuity, leaves the slip free, and selects
a representative of the orbit by a gauge condition (in practice a
small penalty, iteratively refined); see \S\ref{sec:gauge} for when
this is and is not legitimate.

%=====================================================================
\section{The interior slaving}
\label{sec:slaving}

\emph{Status: derived.} Work at one harmonic degree $\ell \ge 1$ on
the radial background, in either dimension. Write the static
(time-independent) equations of the elastic fluid.

\textbf{Step 1 (tangential balance).} The tangential component of the
linearised momentum balance,
$-\nabla p^1 = \rho\,\nabla\phi + \rho^1\nabla\Phi_0$, has no
$\rho^1$-term (gravity is radial), and integrates along the level
surface to
\begin{equation}
  p^1 = -\rho\,\phi
  \qquad (\ell \ge 1).
  \label{eq:p-slaved}
\end{equation}

\textbf{Step 2 (radial balance).} Substituting \eqref{eq:p-slaved}
into the radial component gives
\begin{equation}
  \rho^1 = \frac{1}{g}\frac{\mathrm{d}\rho}{\mathrm{d}r}\,\phi .
  \label{eq:rho-slaved}
\end{equation}
Equations \eqref{eq:p-slaved}--\eqref{eq:rho-slaved} are exactly
Dahlen's interior relations: they follow from momentum balance alone
and hold for \emph{any} static description of the fluid, elastic or
reduced.

\textbf{Step 3 (the elastic description adds kinematics).} The
elastic fluid also carries $\rho^1 = -\Div(\rho\bu)$ and
$\Delta p = -\kappa \Div\bu$ with
$\Delta p = p^1 - \rho g\,u_r$ (using
$\bu\cdot\nabla p^0 = -\rho g\,u_r$). Eliminating $\Div\bu$ between
the two and inserting \eqref{eq:p-slaved}--\eqref{eq:rho-slaved}
yields, after a line of algebra in which the combination
\eqref{eq:N2} appears,
\begin{equation}
  \boxed{\;u_r\,\frac{\rho N^2}{g}
   \;=\; -\,\frac{\rho N^2}{g^2}\,\phi\;}
  \qquad\Longrightarrow\qquad
  u_r = -\frac{\phi}{g}
  \quad\text{whenever } N^2 \neq 0 .
  \label{eq:slaving}
\end{equation}
The radial displacement is \emph{slaved}, pointwise, to the potential
perturbation: fluid particles ride the perturbed equipotentials
(equivalently, the Lagrangian perturbation of the potential-plus-%
position label vanishes). At $N^2 = 0$ the relation degenerates to
$0 = 0$ and $u_r$ is free --- the first appearance of the dichotomy
that organises everything below.

%=====================================================================
\section{Non-existence of the elastic-static solution}
\label{sec:nonexist}

\emph{Status: argued, and displayed numerically.} At the CMB the
kinematics of every formulation that carries a fluid displacement
demands
\begin{equation}
  u_r(\text{fluid side of }\Sigma) = u_r(\text{solid side of }\Sigma),
  \label{eq:matching}
\end{equation}
while \eqref{eq:slaving} assigns the fluid side the value
$-\phi(c)/g(c)$. For a generic load the two assignments differ; the
same trace cannot take two values; hence:

\begin{claim}[no classical static solution]
\label{claim:nonexist}
For $N^2 \ne 0$ and generic loading, the static problem of the
elastic fluid --- welded or slipping alike, since both impose
\eqref{eq:matching} --- has no solution in any function space in
which the fluid displacement has a normal trace at $\Sigma$
(classical, or $H^1$-conforming weak). The over-determination sits at
the interface; no interior rearrangement relieves it.
\end{claim}

Three remarks on what this does and does not say.

\begin{remark}[the variational face]
In weak form the pathology appears as spectrum, not as a visible
contradiction: the stratification contributes a continuum of
near-neutral fluid motions whose energies accumulate at zero
(for the dynamic operator, the essential spectrum is an interval with
endpoint $0$ whenever $N^2 \neq 0$ somewhere). Zero is then not an
isolated eigenvalue with a clean kernel but an accumulation point;
the static operator's range is not closed; and a generic load has a
component outside it. For stable stratification the energy picture is
loss of compactness: minimising sequences push the violation of
\eqref{eq:slaving} into ever-thinner structure at $\Sigma$, the
infimum is approached but not attained, and the limit object leaves
the space. For $N^2 < 0$ the same occurs in saddle form. Turning
this paragraph into a theorem (non-closed range; characterisation of
the compatible loads) is open and worth doing if this material is
published.
\end{remark}

\begin{remark}[what computations produce instead]
A discretised and/or regularised problem is finite-dimensional and
always has an answer. But it approximates an escaping sequence, not a
solution: the computed fluid displacement never settles (it depends
on the mesh and on the regularisation parameter at fixed accuracy),
while gauge-invariant solid observables hover in a band whose size is
controlled by the stratification. The measured bands are in
\S\ref{sec:evidence}. We emphasise: this is not a solver deficiency,
and no amount of resolution removes it.
\end{remark}

\begin{remark}[there is no boundary layer to resolve]
One might expect the solution to consist of the slaved interior plus
a thin matching layer at $\Sigma$. Numerically no such layer appears
at any regularisation we tried: statics contains no mechanism to set
a layer width (no viscosity, no diffusion), so the would-be layer has
zero width --- which is just Claim~\ref{claim:nonexist} again. The
layer becomes real, with a physical width, only when dissipation is
restored (\S\ref{sec:regular}).
\end{remark}

%=====================================================================
\section{The gauge-theoretic statement}
\label{sec:gauge}

\emph{Status: derived (the invariance observation); argued (the
dichotomy).} Two observations place the obstruction relative to the
FRL symmetry.

First, every relabelling of either class satisfies
$\bw\cdot\bn = 0$ on $\Sigma$, so no gauge transformation alters the
normal trace: the incompatibility of
\eqref{eq:slaving}--\eqref{eq:matching} is untouched by any choice of
representative. Second, and more sharply: for $N^2 \ne 0$ the strict
class is tangent to the level surfaces, so its members have
$w_r = 0$ \emph{pointwise} --- the radial displacement of the fluid
is gauge-\emph{invariant} precisely when \eqref{eq:slaving} pins it.
At $N^2 = 0$, where class K makes $u_r$ gauge-variable, the pinning
equation is the one that degenerates. The geometry of the symmetry
group and the solvability of the field equations switch together.

\begin{proposition}[well-posedness on the quotient]
\label{prop:quotient}
The static problem is well posed on the quotient by the fluid
relabellings if and only if $N^2 = 0$. At $N^2 = 0$ the zero modes of
the static operator form a genuine null space (class K), every load
is compatible with it (relabellings are energy- and load-neutral),
and the solution set is a single nonempty orbit; the
Maitra--Al-Attar selection --- normal continuity, free slip, a
representative picked by minimising a gauge functional --- is then a
well-posed choice among existing solutions. For $N^2 \ne 0$ the
degeneracy at zero is essential spectrum rather than kernel, the
solution set is empty for generic loads, and the selection step has
nothing to select from: the gauge-fixed discrete answer it returns is
the regularised object of the previous section, with its band.
\emph{(The ``only if'' direction is Claim~\ref{claim:nonexist}; the
``if'' direction at $N^2=0$ is the class-K lifting argument of the
removability analysis, and is confirmed numerically to $10^{-10}$.)}
\end{proposition}

This is, we believe, the clean way to say what is wrong and what is
right simultaneously: nothing in the relabelling programme is at
fault --- it is exactly the framework in which the dichotomy
kernel-versus-essential-spectrum can be stated --- but its selection
step inherits the solvability of the quotient problem, which is a
property of the stratification, not of the method.

%=====================================================================
\section{Dahlen's method as the secular closure}
\label{sec:dahlen}

\emph{Status: derived (the identification); measured (the
$\kappa$-blindness and the null test).} Dahlen's elimination keeps
only \eqref{eq:p-slaved}--\eqref{eq:rho-slaved} --- which are exact
--- and never introduces a fluid displacement, so
Claim~\ref{claim:nonexist} cannot touch it: its static problem is
well posed (and, in our one-dimensional reference, converges to
machine precision where the elastic-static solves wander). What
happened to the conflict? Reconstruct the displacement implied by
Dahlen's fields through $\rho^1 = -\Div(\rho\bu)$: it is the slaved
field $u_r = -\phi/g$, which does \emph{not} match the solid at
$\Sigma$. The formalism remains consistent because its interface
term
\[
  -\int_{\Sigma} \rho_F\,
     \bigl[\phi\,(\mathbf{m}\cdot\bu') + \phi'\,(\mathbf{m}\cdot\bu)\bigr]\dS
\]
(the ``F3'' term of the implementation; $\mathbf{m}$ the solid's
outward normal, $\rho_F$ the fluid-side density) carries the mass
displaced \emph{across} the boundary as a surface density. In the
displacement picture this is a normal-displacement discontinuity at
the CMB: the fluid slips normally past the boundary and the swept
mass is booked on $\Sigma$. That is the \emph{secular closure}: the
state reached after the fluid has rearranged to hydrostatic
equilibrium in the perturbed potential, however long that takes.

Two sharp consequences, both verified:
(i) at $\ell \ge 1$ the construction never consults the fluid's bulk
modulus, so Dahlen's answer is exactly independent of $\kappa$ in the
core (measured: unchanged to $10^{-11}$ under a thirty-fold scaling);
(ii) at $N^2 = 0$ the reduction coincides with the full elastic
static fluid, so Dahlen, welded and slip must agree exactly there
(measured: $\sim 10^{-10}$ at every degree on \code{aw\_core}). The
well-known failure at $\ell = 0$ --- the elimination discards the
compressional physics, and one constant pressure per fluid region
remains free --- is a separate, understood matter; the practical
remedy (solve $\ell = 0$ with the displacement description) is what
this library and \code{pyslfp} both do.

\paragraph{Aspherical reference states.} The elimination above is a
spherical accident. Its first step --- tangential balance integrating
to $p^1 = -\rho\,\phi$ degree by degree --- used the fact that on a
spherical background the level surfaces are the coordinate spheres
and the harmonics decouple. On a properly aspherical hydrostatic
reference the same balance says only that $p^1 + \rho\,\phi$ is
constant on each level surface of $\Phi_0$: the elimination leaves a
free \emph{function} of $\Phi_0$ per connected fluid region,
coupled to every harmonic through Poisson's equation
(\code{doc/quasi\_static\_models.tex}, assumption M4). The secular
\emph{state} still exists on any geometry --- the fluid rearranged
to hydrostatic equilibrium in the perturbed potential is a
well-defined configuration --- but Dahlen's \emph{route} to it does
not survive asphericity in any practically implementable form. Since
the full-elastic static formulations are the ill-posed ones at
$N^2 \ne 0$, an aspherical model with a non-neutral core currently
has \emph{no} practical well-posed static method at all; the
dissipative route of \S\ref{sec:regular} is then not a consistency
check but the only general-geometry computation of the secular
response on offer. For this library's aspherical programme
(equilibrium figures, aspherical reference states) the point is
practical, not academic.

\begin{remark}[a historical irony]
Dahlen criticised, forcefully, a 1970s proposal of Crossley and
Gubbins that the static deformation problem admits a discontinuity of
radial displacement at the CMB. The analysis above says that his own
static method carries exactly such a discontinuity, implicitly, as
the surface-mass bookkeeping of the secular closure --- and that for
$N^2 \ne 0$ some such discontinuity is \emph{necessary}, since no
static solution without one exists. The only question a formulation
answers is whether it admits the jump openly and is well posed, or
forbids it and loses existence. \emph{(Exact references to be added.)}
\end{remark}

%=====================================================================
\section{What dissipation selects, and routes to regularisation}
\label{sec:regular}

\emph{Status: argued; the Maxwell route is implemented and measured
(\S\ref{sec:maxwell}).}
Statics is the degenerate limit of some evolution, and the
non-existence of \S\ref{sec:nonexist} is the statement that the
limit depends on physics that statics has discarded. Restoring any of
it gives a well-posed problem:

\begin{itemize}
\item \textbf{Fluid viscosity.} With viscosity $\eta$, however small,
  the quasi-static evolution under a steady load is well posed; the
  stress supported by the stratification relaxes as the fluid creeps
  along equipotentials, and the $t \to \infty$ state is the
  rearranged, hydrostatic one --- the Dahlen answer. The interface
  region acquires a genuine layer of $\eta$-dependent width that
  thins as $t$ grows.
\item \textbf{Diffusion of the stratifying quantity.} If the density
  anomaly a displaced parcel carries can diffuse (thermally or
  compositionally), the buoyancy constraint itself erodes;
  weak diffusion again selects the rearranged state. This is, as we
  understand it, Wunsch's point from the 1970s: for a stratified
  fluid the static response is not a property of the inviscid
  adiabatic equations at all --- it is selected by the dissipative
  processes, however weak they are.
\item \textbf{A Maxwell viscoelastic core.} Give the core's shear
  response a Maxwell element (artificial shear modulus $\mu_c$,
  relaxation time $\tau$), apply the load as a step, and relax. At
  every Laplace frequency $s > 0$ the moduli are shifted and the
  problem is uniformly well posed, and the $t \to \infty$ limit is
  the secular state. This is no longer a thought experiment: it is
  the library's production static method, with its own section
  (\S\ref{sec:maxwell}) covering the algorithm, the measured
  trajectory structure (which refines the naive ``relax for a few
  $\tau$'' picture in important ways), the implementation and the
  costs. Its significance is sharpest asphericaly
  (\S\ref{sec:dahlen}): where the Dahlen elimination does not
  survive the geometry, the dissipative route is the only practical
  computation of the secular response, and it works on any mesh this
  library can build.
\item \textbf{What does not work.} A static Tikhonov pin (the small
  deviatoric regulariser $\eps_{\mathrm{reg}}$ of the numerical
  tools) makes each solve unique but selects nothing physical: the
  answer depends on $\eps_{\mathrm{reg}}$ with no limit --- measured,
  and expected, since a static regulariser supplies no arrow of time.
  Likewise the frequency-domain route: the response exists for every
  $\sigma \neq 0$ but has no limit as $\sigma \to 0$ (the limit point
  sits on the essential spectrum); the static limit is genuinely
  non-uniform, which is the spectral restatement of the whole
  phenomenon and connects to the classical secular-constant
  (Longman) discussions. \emph{(References to be added.)}
\end{itemize}

\noindent
The conclusion inverts the naive reading of
Claim~\ref{claim:nonexist}: the elastic-static problem is not
``the'' static problem that Dahlen approximates; Dahlen's closure is
the static limit that every dissipative completion agrees on, and the
elastic-static formulation is the one asking a question without an
answer. What remains true, and matters, is that the \emph{size} of
the gap between any regularised transient and the secular endpoint is
controlled by the stratification, vanishes identically at
$N^2 = 0$, and is invisible in formulations that never resolve the
closure choice explicitly.

%=====================================================================
\section{The Maxwell relaxation method}
\label{sec:maxwell}

\emph{Status: implemented (serial and parallel, every problem class),
measured in one, two and three dimensions; the algorithmic facts
below are measured unless marked otherwise.} The idea
(8 October): since every dissipative completion selects the secular
answer, run one. Make the fluid regions an artificial Maxwell solid,
apply the load as a Heaviside step, and time-step the quasi-static
system by backward Euler until the \emph{solid} region stops moving.
There is then no fluid region in the formulation at all: every step
is a welded solid-elastic solve, the null space is the rigid modes,
and the construction is indifferent to geometry, to the number of
fluid regions, and to their connectivity.

\paragraph{The iteration.} The problem is linear, so the Maxwell
internal variable is the deviatoric fluid strain of a \emph{memory
displacement} $\bw$, and one backward Euler step is
\begin{equation}
  (A + \gamma \hat Q)\,\bu_{n+1} = f + \gamma \hat Q\,\bw_n,
  \qquad
  \bw_{n+1} = \frac{\bw_n + \beta\,\bu_{n+1}}{1+\beta},
  \qquad
  \beta = \frac{\Delta t}{\tau},\quad
  \gamma = \frac{\mu_c}{1+\beta},
  \label{eq:maxwell-step}
\end{equation}
with $A$ the static operator at zero fluid shear and $\hat Q$ the
unit-shear deviatoric fluid stiffness. Three structural facts do all
the work. (i) The fixed point is $\bw = \bu$ with $A\bu = f$: the
relaxed ($\eps \to 0$ welded) solution, \emph{exactly} independent of
$\mu_c$, $\tau$ and the step schedule --- backward Euler's first-order
inaccuracy lives entirely in the transient, which is discarded
(measured: the limit moves by $10^{-13}$ under $\mu_c \times 4$ and
$\Delta t$ halving). (ii) The per-step operator \emph{is} the
gauge-penalty machinery's $A + \eps\mu_c Q$ at
$\eps = 1/(1+\beta)$, and the $\beta \to \infty$ limit of
\eqref{eq:maxwell-step} is precisely the iterated Tikhonov gauge
refinement of \code{doc/gauged\_fluid.md} --- so the long-observed
semi-convergence of those refinements on non-neutral cores was the
Rayleigh--Taylor growth below, seen without knowing it. (iii)
Backward Euler is L-stable, so \emph{escalating} $\beta$ damps every
growing mode: large steps turn the integrator into a continuation
solver ($\gamma \to 0$, toward the static welded solve with a
vanishing deviatoric pin). An A-stable, more accurate scheme
(Crank--Nicolson) would ring instead of damping and has no such
endgame.

\paragraph{The measured trajectory.} Three timescales, not two:
an \emph{elastic} phase contracting at $(1+\beta)^{-1}$ per step
($\sim$10 steps); a \emph{configurational} cascade --- viscous
gravitational relaxation of the artificial fluid, rates
$\lambda/\eta$ with $\lambda$ the weak residual stiffness of
near-equipotential motions, measured down to $10^{-4}\mu_c/\eta$, so
physical times $10^4\,\tau$ and beyond (``a few Maxwell times''
completes only the elastic phase); and, at $N^2 < 0$ only, genuine
Rayleigh--Taylor \emph{growth} of the same viscous-flow type. The
ratio of growth to cascade rates is $\mu_c$-independent, of order
$|N^2|L/g$: on the violently non-neutral \code{fluid\_core}
($\approx 300\times$) there is no plateau window at all and the
finite-time plateau is smeared by exactly the ambiguity band
($3.5\,\%$ at the two-dimensional test resolution), while at
$N^2 \ge 0$ there is no growth phase and the relaxation is clean.
The plateau width is a time-domain display of the band; its
degeneracy is the physics, not the method.

\paragraph{The schedule.} The stopping metric is the relative
per-step \emph{solid} displacement increment only --- precisely the
fields the secular limit defines; the fluid displacement does not
converge at $N^2 \ne 0$ and must not be monitored. The stepper runs
a physical phase at $\beta_0$ (default $3$, which resolves the
plateau --- at $N^2 < 0$ larger $\beta_0$ degrades the returned
band-level answer monotonically, measured to the point of
uselessness at $\beta_0 = 100$; at $N^2 \ge 0$, $\beta_0 = 10$--$30$
is $\sim 3\times$ cheaper and converges with no escalation at all),
detects stagnation of the increment, and then escalates $\beta$
geometrically. Three safeguards, each forced by a measured failure:
\begin{itemize}
\item \textbf{Best-state checkpointing.} The iterate is a function
  of $(\bw, \eps)$ plus a warm start, so the lowest-increment state
  seen is checkpointed by its memory vector alone and restored by
  one warm re-solve. On any terminal event --- iterative-solver
  conditioning floor, escalation cap, step cap --- the best state is
  returned: when the solver cannot complete the continuation, the
  method degrades \emph{to the physical plateau}, the honest
  band-level answer. (Blind escalation instead injected near-kernel
  noise through the conditioning wall and made the answer worse.)
\item \textbf{Increment-validated inexact stepping.} A fixed-point
  iteration is self-correcting, so each step is solved only to a
  fraction (default $0.1$) of the increment it produces, and the
  returned state is polished by one full-tolerance step. The
  validation is essential, not cosmetic: under a loose tolerance a
  warm-started solve can return \emph{unmoved}, faking a zero
  increment and false convergence; each step therefore checks its
  tolerance against the increment it measured and tightens and
  redoes (a warm continuation of the same solve) until the increment
  is resolved.
\item \textbf{An iteration budget} ($20\times$ the first step's
  count) stops escalated solves from burning thousands of iterations
  probing past the conditioning floor.
\end{itemize}

\paragraph{Implementation.} \code{SetMaxwellFluid} on
\code{LinearQuasiStaticProblemBase}: because the per-step operator is
the gauge machinery's, every problem class inherits the mode through
its existing virtual solve --- the referential self-gravitating
class carries its potential block through unchanged --- in one
serial-and-parallel code path. The escalation updates the operator
by scaled sparse addition of a cached unit-penalty matrix (no finite
element reassembly per level). Validated against the exact
gravity-free Lam\'e solutions with a fluid core
(\code{tests/TestFluidGauge*}), and as the seventh architecture of
\code{examples/self\_gravitating\_solvers}.

\paragraph{Cost (measured, degree-2 load).} On the two-dimensional
two-layer disc (order 2, $6474$ unknowns, serial, aw model): the
welded referential solve takes $1738$ MINRES iterations; the
relaxation $9910$ at $\beta_0 = 3$ and $3338$ at $\beta_0 = 30$
(converged, one operator); the slipping architectures $3700$--$4300$
(single-valued) and $17\,000$--$20\,000$ (broken-$\zeta$). On the
three-dimensional two-layer ball (order 1, $7874$ unknowns, 8 MPI
ranks): welded referential $425$ iterations; the relaxation $539$ at
$\beta_0 = 3$ and $\mathbf{314}$ at $\beta_0 = 30$ --- \emph{cheaper
than the welded static solve it replaces}, while handling the fluid
core with no slip machinery, no gauge refinements, and no fluid null
space. On \code{fluid\_core} the solver stops honestly at the
plateau (increment floor $7\times10^{-4}$), landing between the
welded and slip answers --- inside the band, as it must.

\paragraph{The one caveat.} The discrete limit is the \emph{welded}
limit: the fluid-relabelling suppression of the welded space is
inherited, not released (the welded--slip gap at $N^2 \ne 0$ is
gravitational, not elastic, so zero shear does not remove it). At
$\ell \ge 2$ this is inside the band and immaterial; at degree 1 on
non-neutral cores the three-dimensional evidence suggests a
$\sim 2\,\%$ effect riding the soft core--shell translation. The
referee for that question is not the slip formulation (itself a
band-wanderer) but the converged spherical secular answer --- Dahlen
--- and the measurement (relaxation versus Dahlen at degree 1 on
\code{fluid\_core}, with \code{aw\_core} as the null control) is the
one open item between this method and unqualified production status.

%=====================================================================
\section{Numerical evidence}
\label{sec:evidence}

All measurements 7--8 October 2026, with the per-degree radial reference
of the two-dimensional disc family
(\code{benchmarks/disc/disc\_radial.py}; the $\theta$-reduction of
this library's weak forms, solved by one-dimensional $hp$ finite
elements and validated against closed-form solutions of the
gravity-free problems at $10^{-9}$--$10^{-13}$), and with the
finite-element example \code{examples/self\_gravitating\_solvers}
(all six solver architectures of the library on one model and load).
Models: \code{fluid\_core\_2d} (uniform-density core, $N^2 < 0$) and
its mass-matched neutral twin \code{aw\_core\_2d} ($N^2 = 0$ in
closed form); degree-2 loading unless stated.

\begin{center}
\small
\begin{tabular}{@{}llll@{}}
\toprule
measurement & \code{aw\_core\_2d} & \code{fluid\_core\_2d} & reads as \\
\midrule
welded $=$ slip $=$ Dahlen & $\sim 10^{-10}$, every degree &
  spread $\sim 10^{-3}$--$10^{-2}$ & Prop.~\ref{prop:quotient} \\
$p$-refinement of welded/slip & converges, $\sim 10^{-12}$ &
  wanders, $\sim 3\times 10^{-3}$ band & Claim~\ref{claim:nonexist} \\
$p$-refinement of Dahlen & converges & converges ($7$ digits) &
  \S\ref{sec:dahlen} \\
sensitivity to $\eps_{\mathrm{reg}}$ & none ($10^{-8}$--$10^{-12}$) &
  $\sim 10^{-3}$ & no physical selection \\
fluid field tracks \eqref{eq:slaving}? & (n/a: free) & never, at any
  $\eps_{\mathrm{reg}}$ & no-layer remark \\
Dahlen under $\kappa \times 30$ & --- & unchanged to $10^{-11}$ &
  $\kappa$-blindness \\
band about Dahlen vs $N^2\times\frac{1}{30}$ & --- & shrinks
  $\sim 14\times$ (noisy) & band $\to 0$ with $N^2$ \\
FE six-architecture null & all within $1.2\times10^{-3}$ &
  genuine spreads & the FE replica \\
Maxwell route, escalated limit & joins the null ($10^{-10}$) &
  $1.2\times10^{-3}$ from Dahlen & \S\ref{sec:regular} \\
Maxwell route, physical plateau & $=$ limit (no band) &
  $3.5\,\%$ from the limit & the band, in time \\
\bottomrule
\end{tabular}
\end{center}

\medskip\noindent
\textbf{The Maxwell-core run (8 October).} The proposed experiment of
\S\ref{sec:regular}, in the radial tool
(\code{solve\_degree\_maxwell}): backward Euler on the memory
displacement, stopping on the solid increment, step escalation on
stagnation. Measured, degree 2: on \code{aw\_core\_2d} the relaxed
state joins the three-way null at $10^{-10}$ in $28$ steps with no
stagnation; on \code{fluid\_core\_2d} the trajectory has no plateau
window (Rayleigh--Taylor rates $\approx 300\times$ the slowest
configurational ones), the physical plateau sits $3.5\,\%$ off, and
the escalated limit lands on the welded $\eps_{\mathrm{reg}} \to 0$
solution ($5\times10^{-5}$), $1.2\times10^{-3}$ from Dahlen ---
inside the band --- in $25$ solves, exactly independently of $\mu$,
$\tau$ and the step schedule ($10^{-13}$). A third model decides the
$\kappa$-blindness point by pairing: a stably stratified twin
(Adams--Williamson density, $\kappa \times 2$, so $N^2 > 0$)
converges with no growth phase, and on it Dahlen returns the
\emph{identical} value as on \code{aw\_core\_2d} (the M4 elimination
cannot see $\kappa$) while the Maxwell limit moves by
$5\times10^{-4}$: the stratification physics that the secular closure
discards, displayed as a method pair.

\noindent
Fixed-resolution differences \emph{between} the full-elastic methods
on the non-neutral model (welded versus slip, or either versus
Dahlen) lie inside the wandering band and are not resolved
decompositions; only the neutral-twin nulls and Dahlen's convergence
are sharp. Separately, the welded-versus-slip distinction --- the
suppression of non-removable slip by the welded space --- is a
\emph{kinematic} matter orthogonal to the closure question: it too
vanishes at $N^2 = 0$ (class K removes every slip) and is measured as
the $10^{-10}$ null above, but at $N^2 \neq 0$ its magnitude on these
models is likewise inside the band at $\ell \ge 2$ (in three
dimensions it is visible at degree~1, riding the soft core--shell
translation; see \code{doc/planning/open\_issues.md}).

%=====================================================================
\section{Perspective, implications, and open items}
\label{sec:persp}

\paragraph{The status of the formulations.} Where this leaves the
method stack deserves stating plainly, because it inverts the
ranking this project carried into October. Free slip is not physics
as such: it is the $\eta \to 0$ asymptotics of a viscous boundary
layer --- correct about what the \emph{interface} looks like once
the layer has collapsed, and silent about the \emph{interior}
relaxation that collapses it. The measurements agree: at
$N^2 \ne 0$ the slip classes wander in the same resolution band as
welded (freeing the interface kinematics does nothing for the
interior slaving), while Dahlen converges because its closure
contains \emph{both} halves of the dissipative limit, interface and
interior. The hierarchy is therefore: the \textbf{dissipative
completion} (the Maxwell relaxation, \S\ref{sec:maxwell}) is the
formulation-faithful computation, on any geometry; \textbf{Dahlen}
is its converged closed form, available spherically;
\textbf{welded} and \textbf{slip} are finite idealisations of it
--- one representing the layer within a continuous space, the other
replacing it by its zero-width limit --- retained as instruments:
the slip machinery measures the relabelling decomposition (the
welded-suppression question of \S\ref{sec:maxwell}), and the
welded gauged classes are the relaxation's own per-step engine.
At $N^2 = 0$ all of this collapses to a single answer, which is the
standing null test.

\paragraph{Timescales.} The secular regime is entered when the
loading timescale exceeds the buoyancy timescale. Even a core
stratification as feeble as $N^2 \sim 10^{-15}\,\mathrm{s}^{-2}$ ---
six orders below commonly discussed outer-core values --- gives
buoyancy periods of order years, against GIA loading periods of
millennia: $\sigma/N \lesssim 10^{-3}$ in the most unfavourable
reading, and $\sim 10^{-7}$ for plausible values. For GIA, therefore,
the Dahlen closure is simply correct, the elastic-static band is
transient physics the Earth has long since relaxed through, and the
entire effect is negligible against (for example) three-dimensional
mantle-viscosity variations. The practical force of this note for GIA
is indirect: it explains the non-neutral resolution floor seen in
benchmark ladders, disciplines what comparisons against Dahlen-type
references (such as \code{pyslfp}) can and cannot resolve, and
identifies the neutral twin \code{aw\_core} as the configuration on
which everything is clean.

\paragraph{Where it is not academic.} Observables or bodies for which
core stratification is itself the target --- stably stratified layers
at the top of the core, other planets, or any inversion in which
$N^2$ is a parameter --- inherit an $N^2$-controlled modelling choice
that is invisible in standard treatments because the closure is
implicit. There the honest statement is: the static response is
defined only up to the secular closure, Dahlen computes that closure,
and the distance to any finite-time (partially relaxed) state is a
transient that requires the dissipative physics to evaluate.

\paragraph{Relation to the dynamic ill-posedness argument (DAMTP
2025 talk).} A sketched argument of the author's (slides,
\code{\~{}/Documents/Talks/DAMTP2025}; to be completed) holds that
the \emph{dynamic} equations in a non-neutrally stratified fluid are
ill-posed through a loss of regularity: writing the fluid
elasto-dynamic operator as $A = C^*(C+D+E) + FC$, with $C$ the
divergence-type operator and $F$ the buoyancy coupling that
\emph{vanishes exactly in the neutral case}, the kernel sector of
the resolvent reduction is driven by the acoustic sector through the
composite $FC$ --- energy transfer from acoustic to internal modes.
The essential mechanism is a \emph{loss of regularity}: because of
the factor of $C$ the coupling is an operator of order one, so the
internal sector is forced by a quantity one derivative rougher than
its space, and the $\ker C$ dynamics (free, with no smoothing)
cannot pay the derivative back --- $u_0$ does not remain in the
right space, at any frequency. The explicit $p^{-1}$ factor is the
additional low-frequency amplification on top of this. The present
note is the $p \to 0$ endpoint of the same structure: at $F = 0$
the kernel sector decouples as pure gauge and the static quotient
problem is well posed (Proposition~\ref{prop:quotient}); at
$F \neq 0$ the driving has no static limit, which is
Claim~\ref{claim:nonexist},
the non-uniform $\sigma \to 0$ limit of \S\ref{sec:regular}, and
the measured band, while the interfacial loss of compactness is the
spatial face of the dynamic loss of regularity. The two gradations
are consistent: dynamically, mild solutions survive while
derivatives are progressively lost; statically, where there is no
time left to lose them in, existence itself fails. In particular the open theorem below and the talk's
to-be-completed argument are one piece of mathematics --- the
essential-spectrum characterisation there reduces to the Fredholm
property of $1_Z - p^{-1}D\,(QB|_{\ker C} + p\,1)^{-1}QF$, whose
$p \to 0$ corner is exactly the static solvability question ---
and the measurements of \S\ref{sec:evidence} are numerical evidence
for the dynamic conjecture as well.

\paragraph{Open items.}
(i) The functional-analytic theorem behind
Claim~\ref{claim:nonexist} and Proposition~\ref{prop:quotient} ---
now best approached jointly with the dynamic argument above.
(ii) The Maxwell relaxation (\S\ref{sec:maxwell}): the degree-1
adjudication (relaxation versus Dahlen on \code{fluid\_core},
\code{aw\_core} as control --- the one open question between the
method and unqualified production status), the benchmark-grade
Love-number runs, and the aspherical demonstration --- a static
fluid-core response on a geometry where no other method exists.
(iii) A clean measurement of the band's $N^2$-scaling (the
one-dimensional tool does each point in seconds; the current
two-resolution estimator is noisy).
(iv) The three-dimensional \code{aw\_core} runs: the degree-1
welded--slip null test, the Dahlen-path $h$-ladder (deciding whether
the non-neutral floor explains the outstanding
\code{fluid\_core} ladder residuals), and the decomposition of the
\code{pyslfp} offsets.
(v) References: Dahlen's 1970s static-deformation papers, the
Crossley--Gubbins exchange, Longman's problem, and Wunsch's
dissipation argument, to be added precisely.

\end{document}
